\documentclass[12pt, a4paper]{article}

\usepackage[utf8]{inputenc}
\usepackage[T1]{fontenc}
\usepackage{amsmath, amssymb, amsthm, mathtools}
\usepackage{geometry}
\usepackage{hyperref}

\geometry{margin=1in}
\hypersetup{colorlinks=true, linkcolor=blue, citecolor=blue, urlcolor=blue}

\newtheorem{theorem}{Theorem}[section]
\newtheorem{lemma}[theorem]{Lemma}
\newtheorem{proposition}[theorem]{Proposition}
\newtheorem{corollary}[theorem]{Corollary}
\theoremstyle{definition}
\newtheorem{definition}[theorem]{Definition}
\newtheorem{remark}[theorem]{Remark}

\newcommand{\R}{\mathbb{R}}
\newcommand{\N}{\mathbb{N}}
\newcommand{\supp}{\operatorname{supp}}
\newcommand{\tv}{\operatorname{tv}}
\newcommand{\rev}{\operatorname{rev}}
\newcommand{\conv}{\operatorname{conv}}

\title{%
  Proving Lower Bounds on the Autoconvolution Constant $C_{1a}$\\[6pt]
  \large Via the Cloninger--Steinerberger Branch-and-Prune Algorithm
}
\author{Jai Bajaj}
\date{\today}

\begin{document}
\maketitle

\begin{abstract}
We outline a proof framework for establishing lower bounds on the
autoconvolution constant $C_{1a}$ via exhaustive branch-and-prune
computation.  TODO: expand.
\end{abstract}

\tableofcontents
\newpage

% ======================================================================
\section{Introduction}\label{sec:intro}
% ======================================================================

Consider the class of nonnegative, integrable functions supported on
the interval $[-\tfrac14, \tfrac14]$:
\[
  \mathcal{F} \;=\; \Bigl\{
    f : \R \to \R_{\ge 0}
    \;\Big|\;
    \supp(f) \subseteq \bigl[-\tfrac14, \tfrac14\bigr],\;
    \int f > 0
  \Bigr\}.
\]

The \emph{autoconvolution} of $f$ is
\[
  (f * f)(t) \;=\; \int_{-\infty}^{\infty} f(x)\, f(t - x)\, dx.
\]
Since $\supp(f) \subseteq [-\tfrac14, \tfrac14]$, the
autoconvolution is supported on $[-\tfrac12, \tfrac12]$.

\begin{definition}
The \emph{autoconvolution constant} is
\[
  C_{1a} \;=\; \inf_{f \in \mathcal{F}} \;
  \frac{\displaystyle\max_{|t| \le 1/2}\, (f*f)(t)}%
       {\displaystyle\left(\int f\right)^{\!2}}.
\]
\end{definition}

By homogeneity, we may normalize $\int f = 1$, so
$C_{1a} = \inf_{f} \|f*f\|_{L^\infty}$.

The best known bounds are $1.2802 \leq C_{1a} \leq 1.5029$.  The
upper bound comes from exhibiting explicit functions; the lower bound
(Cloninger \& Steinerberger, 2017) requires showing that \emph{no}
function can achieve a peak below the threshold.  This is the subject
of our proof.

The proof strategy:
\begin{enumerate}
  \item Reduce from continuous functions to finite-dimensional bin averages.
  \item Discretize the simplex to a finite lattice, with error bounds.
  \item Exhaustively check every lattice point, using pruning and a
    multiscale cascade to make this tractable.
\end{enumerate}


% ======================================================================
\section{From Functions to Bin Averages}\label{sec:bin}
% ======================================================================

% TODO: partition, bin averages, Lemma 1 (windowed lower bound)


% ======================================================================
\section{Discretization of the Simplex}\label{sec:disc}
% ======================================================================

% TODO: lattice B_{n,m}, 1/m-net, correction 2/m + 1/m^2, dynamic correction


% ======================================================================
\section{The Windowed Test Value}\label{sec:testvalue}
% ======================================================================

% TODO: definition, prefix sums, reversal symmetry


% ======================================================================
\section{Pruning Strategies}\label{sec:pruning}
% ======================================================================

% TODO: asymmetry, block-sum, two-max, canonical symmetry


% ======================================================================
\section{The Multiscale Cascade}\label{sec:cascade}
% ======================================================================

% TODO: dyadic splitting, energy cap, cascade soundness, canonical handling


% ======================================================================
\section{Correctness and Verification}\label{sec:correctness}
% ======================================================================

% TODO: four invariants, computational checks


% ======================================================================
\section{Discussion}\label{sec:discussion}
% ======================================================================

% TODO: relation to CS17, limitations


\begin{thebibliography}{9}

\bibitem{CS17}
A.~Cloninger and S.~Steinerberger,
\emph{On Suprema of Autoconvolutions with an Application to Sidon Sets},
Proc.\ AMS \textbf{145}(8), 3191--3200, 2017.
\href{https://arxiv.org/abs/1403.7988}{arXiv:1403.7988}.

\bibitem{MV10}
M.~Matolcsi and C.~Vinuesa,
\emph{Improved bounds on the supremum of autoconvolutions},
J.\ Math.\ Anal.\ Appl.\ \textbf{372}(2), 439--447, 2010.

\end{thebibliography}

\end{document}
